\documentclass[11pt,a4paper]{article}
\usepackage[utf8]{inputenc}
\usepackage[spanish]{babel}
\usepackage{amsmath,amssymb,amsfonts}
\usepackage{geometry}
\geometry{top=2.5cm, bottom=2.5cm, left=2.5cm, right=2.5cm}
\usepackage{listings}
\usepackage{xcolor}
\usepackage{hyperref}
\usepackage{titlesec}
\usepackage{tikz}
\usepackage{xspace}
\usepackage{graphicx}
\usepackage{fancyhdr}
\usetikzlibrary{shapes.geometric, arrows.meta, positioning, calc}

% Configuración de colores para código
\definecolor{codegreen}{rgb}{0,0.6,0}
\definecolor{codegray}{rgb}{0.5,0.5,0.5}
\definecolor{codepurple}{rgb}{0.58,0.82,0.82}
\definecolor{backcolour}{rgb}{0.96,0.96,0.96}
\definecolor{keywordcolor}{rgb}{0.13,0.13,0.75}
\definecolor{numbercolor}{rgb}{0.8,0.4,0}

\lstdefinestyle{luastyle}{
    backgroundcolor=\color{backcolour},   
    commentstyle=\color{codegreen},
    keywordstyle=\color{keywordcolor}\bfseries,
    numberstyle=\tiny\color{codegray},
    stringstyle=\color{red!70!black},
    basicstyle=\ttfamily\small,
    breakatwhitespace=false,         
    breaklines=true,                 
    captionpos=b,                    
    keepspaces=true,                 
    numbers=left,                    
    numbersep=5pt,                  
    showspaces=false,                
    showstringspaces=false,
    showtabs=false,                  
    tabsize=2,
    frame=single,
    extendedchars=true,
    inputencoding=utf8,
    literate=
        {á}{{\'a}}1 {é}{{\'e}}1 {í}{{\'i}}1 {ó}{{\'o}}1 {ú}{{\'u}}1
        {Á}{{\'A}}1 {É}{{\'E}}1 {Í}{{\'I}}1 {Ó}{{\'O}}1 {Ú}{{\'U}}1
        {ñ}{{\~n}}1 {Ñ}{{\~N}}1 {¿}{{\?`}}1 {¡}{{!`}}1
        {0}{{{\color{numbercolor}0}}}1
        {1}{{{\color{numbercolor}1}}}1
        {2}{{{\color{numbercolor}2}}}1
        {3}{{{\color{numbercolor}3}}}1
        {4}{{{\color{numbercolor}4}}}1
        {5}{{{\color{numbercolor}5}}}1
        {6}{{{\color{numbercolor}6}}}1
        {7}{{{\color{numbercolor}7}}}1
        {8}{{{\color{numbercolor}8}}}1
        {9}{{{\color{numbercolor}9}}}1
}

\lstdefinelanguage{lua}{
  morekeywords={local, function, end, if, then, else, elseif, while, do, repeat, until, for, in, return, break, nil, true, false, assert, error, import, tostring, print},
  sensitive=true,
  morecomment=[l]{--},
  morecomment=[s][\color{codegreen}]{--[===[}{]===]},
  morecomment=[s][\color{codegreen}]{--[==[}{]==]},
  morecomment=[s][\color{codegreen}]{--[=[}{]=]},
  morecomment=[s][\color{codegreen}]{--[[}{]]},
  morestring=[b]",
  morestring=[b]',
  morecomment=[s][\color{red!70!black}]{[===[}{]===]},
  morecomment=[s][\color{red!70!black}]{[==[}{]==]},
  morecomment=[s][\color{red!70!black}]{[=[}{]=]},
  morecomment=[s][\color{red!70!black}]{[[}{]]}
}

\lstset{style=luastyle}

% --- ENCABEZADOS Y PIE DE PÁGINA ---
\pagestyle{fancy}
\fancyhf{} 
\rhead{\includegraphics[height=1.2cm]{logo.png}}
\lhead{\small \textbf{Base de Datos} -- Actividad 3}
\rfoot{Página \thepage}

% --- DATOS DE LA PORTADA ---
\title{\textbf{Administración y Diseño de Base de Datos: Concesionaria de Automóviles}\\ \large Análisis Operativo, Modelo Relacional, Normalización y Carga Automatizada en LuaJIT}
\author{\textbf{Nombre:} GARCIA CAPISTRAN JOSE CARLOS \\ \textbf{Matrícula:} AL07172840 \\ \textbf{Materia:} Base de Datos}
\date{\today}

\begin{document}
\maketitle
\thispagestyle{fancy}

\hrule
\vspace{0.5cm}

\section{Análisis de la Industria y Organización de Información}

El modelo de negocio de una red de \textbf{Concesionarios de Automóviles} administra el inventario multimarca, la asignación de vendedores a salas de exhibición, el registro de unidades físicamente disponibles y el procesamiento transaccional de ventas tanto de contado como a crédito.

\subsection{Giro de la Empresa y Uso de la Información}
La empresa pertenece al giro automotriz y de servicios comerciales. Maneja catálogos de modelos con especificaciones técnicas (potencia, tipo de motor, cilindraje), números de serie de vehículos en stock, datos de contacto de ejecutivos de venta y transacciones comerciales estructuradas.

\subsection{Almacenamiento Tradicional vs. Base de Datos Relacional}
Tradicionalmente, las concesionarias organizan su información en hojas de cálculo no integradas compartidas por red local o correo electrónico:
\begin{itemize}
    \item \textbf{Excel de Catálogo e Inventario}: Registro de modelos, números de serie, precios base y ubicaciones (local, almacén, servicio oficial).
    \item \textbf{Excel de Vendedores y Comisiones}: Directorios con teléfonos, WhatsApp y asignación a concesionarias.
    \item \textbf{Excel de Registro de Ventas}: Mapeo de clientes, matrículas asignadas, fechas de entrega y modalidades de pago.
\end{itemize}

\textbf{Deficiencias del modelo tradicional}: Redundancia de datos, riesgo de duplicidad en números de serie (VIN), falta de integridad referencial al eliminar modelos descontinuados y desactualización del inventario real ante ventas simultáneas.

\section{Diseño Conceptual y Restricciones de Integridad}

Se diseñó un modelo conceptual centralizado con las siguientes entidades principales y sus atributos clave:

\begin{itemize}
    \item \textbf{Concesionarios}: \texttt{id\_concesionario}, \texttt{nombre}, \texttt{domicilio}, \texttt{telefono}.
    \item \textbf{Vendedores}: \texttt{id\_vendedor}, \texttt{id\_concesionario}, \texttt{nombre}, \texttt{domicilio}, \texttt{telefono}, \texttt{whatsapp}.
    \item \textbf{Modelos\_Autos}: \texttt{id\_modelo}, \texttt{marca}, \texttt{modelo}, \texttt{anho}, \texttt{precio\_base}, \texttt{descuento}, \texttt{tipo\_motor}, \texttt{potencia}, \texttt{cilindros}.
    \item \textbf{Automoviles\_Stock}: \texttt{numero\_serie}, \texttt{id\_modelo}, \texttt{id\_concesionario}, \texttt{ubicacion}.
    \item \textbf{Ventas}: \texttt{id\_venta}, \texttt{numero\_serie}, \texttt{id\_vendedor}, \texttt{precio\_venta}, \texttt{modo\_pago}, \texttt{fecha\_entrega}, \texttt{matricula}, \texttt{origen\_adquisicion}.
    \item \textbf{Ventas\_Acumuladas}: \texttt{id\_simulacion}, \texttt{fecha}, \texttt{modelo}, \texttt{vendedor}, \texttt{total}.
\end{itemize}

\subsection{Restricciones de Integridad}
\begin{enumerate}
    \item \textbf{Integridad de Entidad}: Claves primarias únicas definidas en cada tabla (\texttt{id\_concesionario}, \texttt{id\_vendedor}, \texttt{id\_modelo}, \texttt{numero\_serie}, \texttt{id\_venta}, \texttt{id\_simulacion}).
    \item \textbf{Integridad Referencial}: Llaves foráneas con reglas \texttt{ON DELETE CASCADE / RESTRICT} y \texttt{ON UPDATE CASCADE} para preservar la consistencia entre stock, vendedores y ventas.
    \item \textbf{Integridad de Dominio}: Campos restrictivos con enumeraciones (\texttt{ENUM('Contado', 'Credito')}, \texttt{ENUM('Stock', 'Fabrica')}) y restricciones de completitud (\texttt{NOT NULL}).
\end{enumerate}

\newpage

\section{Diagrama Entidad-Relación}

A continuación se presenta la arquitectura conceptual y las relaciones cardinales del sistema de concesionarios:

\begin{figure}[h!]
\centering
\resizebox{0.95\textwidth}{!}{
    \begin{tikzpicture}[
        node distance=1.6cm and 2.2cm,
        entity/.style={rectangle, draw=blue!80, fill=blue!5, thick, minimum width=2.8cm, minimum height=1cm, align=center, font=\small},
        rel/.style={diamond, draw=red!80, fill=red!5, thick, aspect=1.8, inner sep=1pt, align=center, font=\scriptsize},
        card/.style={font=\tiny\bfseries, color=gray!90!black}
    ]
        % Entidades
        \node[entity] (concesionarios) {\textbf{concesionarios}\\ \tiny id\_concesionario (PK)};
        \node[entity, right=3.0cm of concesionarios] (vendedores) {\textbf{vendedores}\\ \tiny id\_vendedor (PK)};
        \node[entity, below=2.5cm of concesionarios] (stock) {\textbf{automoviles\_stock}\\ \tiny numero\_serie (PK)};
        \node[entity, left=3.0cm of stock] (modelos) {\textbf{modelos\_autos}\\ \tiny id\_modelo (PK)};
        \node[entity, right=3.0cm of stock] (ventas) {\textbf{ventas}\\ \tiny id\_venta (PK)};

        % Relaciones
        \node[rel] (r_emplea) at ($(concesionarios)!0.5!(vendedores)$) {Emplea};
        \node[rel] (r_almacena) at ($(concesionarios)!0.5!(stock)$) {Almacena};
        \node[rel] (r_clasifica) at ($(modelos)!0.5!(stock)$) {Clasifica};
        \node[rel] (r_realiza) at ($(vendedores)!0.5!(ventas)$) {Realiza};
        \node[rel] (r_vende) at ($(stock)!0.5!(ventas)$) {Suministra};

        % Conexiones
        \draw[thick] (concesionarios) -- (r_emplea) -- (vendedores);
        \draw[thick] (concesionarios) -- (r_almacena) -- (stock);
        \draw[thick] (modelos) -- (r_clasifica) -- (stock);
        \draw[thick] (vendedores) -- (r_realiza) -- (ventas);
        \draw[thick] (stock) -- (r_vende) -- (ventas);
    \end{tikzpicture}
}
\caption{Diagrama Entidad-Relación del Sistema de Concesionarios.}
\label{fig:er_concesionario}
\end{figure}

\newpage
\section{Estructura de Datos y Diagrama Relacional}

Representación detallada de las tablas creadas en MariaDB InnoDB:

\begin{figure}[h!]
\centering
\resizebox{0.95\textwidth}{!}{
    \begin{tikzpicture}[
        tablebox/.style={
            draw=black!70,
            fill=white,
            thick,
            rounded corners=2pt,
            inner sep=0pt,
            font=\small
        }
    ]
        % FILA 1
        \node[tablebox] (concesionarios) {
            \begin{tabular}{l|r}
                \multicolumn{2}{c}{\textbf{concesionarios}} \\ \hline
                \textcolor{orange!80!black}{$\mathbf{\mu}$}\ id\_concesionario & \textit{INT (PK)} \\
                nombre & \textit{VARCHAR(100)} \\
                domicilio & \textit{VARCHAR(200)} \\
                telefono & \textit{VARCHAR(20)} \\
            \end{tabular}
        };

        \node[tablebox, right=2.0cm of concesionarios] (vendedores) {
            \begin{tabular}{l|r}
                \multicolumn{2}{c}{\textbf{vendedores}} \\ \hline
                \textcolor{orange!80!black}{$\mathbf{\mu}$}\ id\_vendedor & \textit{INT (PK)} \\
                id\_concesionario & \textit{INT (FK)} \\
                nombre & \textit{VARCHAR(100)} \\
                telefono & \textit{VARCHAR(20)} \\
                whatsapp & \textit{VARCHAR(20)} \\
            \end{tabular}
        };

        % FILA 2
        \node[tablebox, below=2.5cm of concesionarios] (modelos) {
            \begin{tabular}{l|r}
                \multicolumn{2}{c}{\textbf{modelos\_autos}} \\ \hline
                \textcolor{orange!80!black}{$\mathbf{\mu}$}\ id\_modelo & \textit{INT (PK)} \\
                marca & \textit{VARCHAR(50)} \\
                modelo & \textit{VARCHAR(50)} \\
                anho & \textit{INT} \\
                precio\_base & \textit{DECIMAL(12,2)} \\
            \end{tabular}
        };

        \node[tablebox, right=2.0cm of modelos] (stock) {
            \begin{tabular}{l|r}
                \multicolumn{2}{c}{\textbf{automoviles\_stock}} \\ \hline
                \textcolor{orange!80!black}{$\mathbf{\mu}$}\ numero\_serie & \textit{VARCHAR(50) (PK)} \\
                id\_modelo & \textit{INT (FK)} \\
                id\_concesionario & \textit{INT (FK)} \\
                ubicacion & \textit{VARCHAR(50)} \\
            \end{tabular}
        };

        \node[tablebox, right=2.0cm of stock] (ventas) {
            \begin{tabular}{l|r}
                \multicolumn{2}{c}{\textbf{ventas}} \\ \hline
                \textcolor{orange!80!black}{$\mathbf{\mu}$}\ id\_venta & \textit{INT (PK)} \\
                numero\_serie & \textit{VARCHAR(50) (FK)} \\
                id\_vendedor & \textit{INT (FK)} \\
                precio\_venta & \textit{DECIMAL(12,2)} \\
                modo\_pago & \textit{ENUM} \\
            \end{tabular}
        };

        % Conexiones
        \draw[thick] (concesionarios.east) -- (vendedores.west);
        \draw[thick] (concesionarios.south) -- (stock.north);
        \draw[thick] (modelos.east) -- (stock.west);
        \draw[thick] (stock.east) -- (ventas.west);
        \draw[thick] (vendedores.south) |- (ventas.north);
    \end{tikzpicture}
}
\caption{Modelo Relacional de Concesionarias y Ventas.}
\end{figure}

\subsection{Procedimientos para Pruebas de Integridad y Conectividad}
\begin{itemize}
    \item \textbf{Prueba de Conectividad en LuaJIT}: Inicialización exitosa del socket de MariaDB mediante el módulo \texttt{cmariadb.so}.
    \item \textbf{Prueba de Eliminación en Cascada/Restricción}: Verificación de la desactivación temporal del control de claves foráneas con \texttt{SET FOREIGN\_KEY\_CHECKS = 0;} durante ejecuciones masivas.
\end{itemize}

\section{Diseño Lógico y Proceso de Normalización}

Se aplicaron rigurosamente las primeras tres formas normales (1NF, 2NF, 3NF):

\begin{enumerate}
    \item \textbf{Primera Forma Normal (1NF)}: Todos los atributos son atómicos. Se eliminaron listas de automóviles o teléfonos múltiples concatenados en una columna.
    \item \textbf{Segunda Forma Normal (2FN)}: Los atributos no clave de cada tabla dependen completamente de su clave primaria. Los detalles del modelo de auto se separaron del registro individual del número de serie en stock.
    \item \textbf{Tercera Forma Normal (3FN)}: Eliminación de dependencias transitivas. Los datos del concesionario o vendedor no se almacenan directamente en la tabla de ventas, sino a través de claves foráneas referenciadas.
\end{enumerate}

\newpage

\section{Configuración de Accesos y Privilegios}

Configuración de usuarios en MariaDB para acceso mediante socket o TCP local:

\begin{lstlisting}[language=sql, caption={Configuración del usuario de aplicación}]
CREATE USER IF NOT EXISTS 'lua_client'@'localhost' IDENTIFIED BY '12345';
GRANT ALL PRIVILEGES ON autos_concesionario_db.* TO 'lua_client'@'localhost';
FLUSH PRIVILEGES;
\end{lstlisting}

\section{Implementación del Script en LuaJIT (program.actividad3.lua)}

El siguiente script ejecuta la creación del esquema DDL, la siembra automatizada de 20 registros por tabla sin valores nulos, y las consultas de actualización, eliminación y filtrado:

\begin{lstlisting}[language=lua, caption={Script en LuaJIT para automatización de la Base de Datos (program.actividad3.lua)}]
import("cmariadb", "system")

-- 1. Conexión a la base de datos
local db, err = cmariadb.connect({
    host = "127.0.0.1",
    user = "lua_client",
    password = "12345",
    client_mode = cmariadb.CLIENT_MODE.MULTIPLE_STATEMENTS
})
if not db then
    error("[ERROR] No se pudo conectar a MariaDB: " .. tostring(err))
end
print("[INFO] Conexión establecida con MariaDB.")

-- Creación y selección de la base de datos
assert(db:query("DROP DATABASE IF EXISTS autos_concesionario_db;"))
assert(db:query("CREATE DATABASE IF NOT EXISTS autos_concesionario_db CHARACTER SET utf8mb4 COLLATE utf8mb4_unicode_ci;"))
assert(db:query("USE autos_concesionario_db;"))
print("[INFO] Base de datos 'autos_concesionario_db' creada y seleccionada.")

-- 2. Definición del Esquema DDL (Relaciones, Llaves Primarias y Foráneas)
local ddl_schema = [[
    CREATE TABLE IF NOT EXISTS concesionarios (
        id_concesionario INT AUTO_INCREMENT PRIMARY KEY,
        nombre VARCHAR(100) NOT NULL,
        domicilio VARCHAR(200) NOT NULL,
        telefono VARCHAR(20) NOT NULL
    ) ENGINE=InnoDB;

    CREATE TABLE IF NOT EXISTS vendedores (
        id_vendedor INT AUTO_INCREMENT PRIMARY KEY,
        id_concesionario INT NOT NULL,
        nombre VARCHAR(100) NOT NULL,
        domicilio VARCHAR(200) NOT NULL,
        telefono VARCHAR(20) NOT NULL,
        whatsapp VARCHAR(20) NOT NULL,
        FOREIGN KEY (id_concesionario) REFERENCES concesionarios(id_concesionario) ON DELETE CASCADE ON UPDATE CASCADE
    ) ENGINE=InnoDB;

    CREATE TABLE IF NOT EXISTS modelos_autos (
        id_modelo INT AUTO_INCREMENT PRIMARY KEY,
        marca VARCHAR(50) NOT NULL,
        modelo VARCHAR(50) NOT NULL,
        anho INT NOT NULL,
        precio_base DECIMAL(12,2) NOT NULL,
        descuento DECIMAL(5,2) DEFAULT 0.00,
        tipo_motor VARCHAR(50) NOT NULL,
        potencia INT NOT NULL,
        cilindros INT NOT NULL
    ) ENGINE=InnoDB;

    CREATE TABLE IF NOT EXISTS automoviles_stock (
        numero_serie VARCHAR(50) PRIMARY KEY,
        id_modelo INT NOT NULL,
        id_concesionario INT NOT NULL,
        ubicacion VARCHAR(50) NOT NULL, -- Local, Sucursal, Almacen, Servicio Oficial
        FOREIGN KEY (id_modelo) REFERENCES modelos_autos(id_modelo) ON DELETE CASCADE ON UPDATE CASCADE,
        FOREIGN KEY (id_concesionario) REFERENCES concesionarios(id_concesionario) ON DELETE CASCADE ON UPDATE CASCADE
    ) ENGINE=InnoDB;

    CREATE TABLE IF NOT EXISTS ventas (
        id_venta INT AUTO_INCREMENT PRIMARY KEY,
        numero_serie VARCHAR(50) NOT NULL,
        id_vendedor INT NOT NULL,
        precio_venta DECIMAL(12,2) NOT NULL,
        modo_pago ENUM('Contado', 'Credito') NOT NULL,
        fecha_entrega DATE NOT NULL,
        matricula VARCHAR(20) NOT NULL,
        origen_adquisicion ENUM('Stock', 'Fabrica') NOT NULL,
        FOREIGN KEY (numero_serie) REFERENCES automoviles_stock(numero_serie) ON DELETE RESTRICT ON UPDATE CASCADE,
        FOREIGN KEY (id_vendedor) REFERENCES vendedores(id_vendedor) ON DELETE RESTRICT ON UPDATE CASCADE
    ) ENGINE=InnoDB;

    CREATE TABLE IF NOT EXISTS ventas_acumuladas (
        id_simulacion INT AUTO_INCREMENT PRIMARY KEY,
        fecha DATE NOT NULL,
        modelo VARCHAR(100) NOT NULL,
        vendedor VARCHAR(100) NOT NULL,
        total DECIMAL(12,2) NOT NULL
    ) ENGINE=InnoDB;
]]

assert(db:multi_query(ddl_schema))
print("[INFO] Estructura relacional creada exitosamente.")

-- 3. Inserción automatizada de 20 registros NOT NULL por tabla
math.randomseed(os.time())

local marcas_arr = {"Audi", "Suzuki", "Honda", "Volkswagen", "BMW", "Toyota", "Nissan", "Mazda"}
local nombres_arr = {"Carlos", "Ana", "Luis", "Maria", "Jorge", "Sofia", "Diego", "Valeria"}
local apellidos_arr = {"Gomez", "Lopez", "Perez", "Martines", "Sanchez", "Ramirez"}

local function rand_elt(t) return t[math.random(#t)] end

-- Generar datos en bloques
local con_vals, vend_vals, mod_vals, auto_vals, vent_vals, sim_vals = {}, {}, {}, {}, {}, {}

for i = 1, 20, 1 do
    -- Concesionarios
    table.insert(con_vals, string.format("('Concesionario Central %d', 'Av. Reforma %d', '555000%02d')", i, i*10, i))
    
    -- Vendedores (asociados a concesionarios 1 a 20)
    table.insert(vend_vals, string.format("(%d, '%s %s', 'Calle Sur %d', '555111%02d', '555111%02d')", 
        ((i-1)%5)+1, rand_elt(nombres_arr), rand_elt(apellidos_arr), i, i, i))
    
    -- Modelos de Autos (mezclando años antes y después de 2020 para la prueba de DELETE)
    local anho_val = (i <= 5) and (2017 + (i % 3)) or 2021 + (i % 5)
    table.insert(mod_vals, string.format("('%s', 'Modelo-%d', %d, %d.00, %d.00, 'Turbo', %d, %d)", 
        rand_elt(marcas_arr), i, anho_val, 140000 + (i * 12000), (i%2==0) and 5.00 or 0.00, 150 + i*10, (i%2==0)?4:6))
    
    -- Automóviles en Stock
    table.insert(auto_vals, string.format("('SN-VIN-%04d', %d, %d, '%s')", 
        i, i, ((i-1)%5)+1, (i%2==0) and "Local Ventas" or "Servicio Oficial"))
        
    -- Ventas
    table.insert(vent_vals, string.format("('SN-VIN-%04d', %d, %d.00, '%s', '2026-08-%02d', 'ABC-%03d', '%s')", 
        i, ((i-1)%5)+1, 150000.00 + (i * 5000), (i%2==0) and "Contado" or "Credito", (i%28)+1, i, (i%3==0) and "Fabrica" or "Stock"))
end

-- 25 Simulaciones para ventas_acumuladas
for i = 1, 25, 1 do
    table.insert(sim_vals, string.format("('2026-08-%02d', 'Modelo-%d', 'Vendedor %d', %d.00)", 
        (i%28)+1, (i%20)+1, (i%5)+1, 140000 + (i * 8500)))
end

local dml_inserts = string.format([[
    SET FOREIGN_KEY_CHECKS = 0;
    INSERT INTO concesionarios (nombre, domicilio, telefono) VALUES %s;
    INSERT INTO vendedores (id_concesionario, nombre, domicilio, telefono, whatsapp) VALUES %s;
    INSERT INTO modelos_autos (marca, modelo, anho, precio_base, descuento, tipo_motor, potencia, cilindros) VALUES %s;
    INSERT INTO automoviles_stock (numero_serie, id_modelo, id_concesionario, ubicacion) VALUES %s;
    INSERT INTO ventas (numero_serie, id_vendedor, precio_venta, modo_pago, fecha_entrega, matricula, origen_adquisicion) VALUES %s;
    INSERT INTO ventas_acumuladas (fecha, modelo, vendedor, total) VALUES %s;
    SET FOREIGN_KEY_CHECKS = 1;
]], table.concat(con_vals, ","), table.concat(vend_vals, ","), table.concat(mod_vals, ","), 
    table.concat(auto_vals, ","), table.concat(vent_vals, ","), table.concat(sim_vals, ","))

assert(db:multi_query(dml_inserts))
print("[INFO] Registros semilla insertados correctamente en todas las tablas.")

-- ==========================================
-- 4. EJECUCIÓN DE CONSULTAS Y OPERACIONES SOLICITADAS
-- ==========================================

print("\n--- A. VENTAS DE UNA SOLA CONCESIONARIA (LIMIT 5) ---")
local q_limit = db:query([[
    SELECT v.id_venta, c.nombre AS concesionario, v.precio_venta, v.modo_pago, v.matricula 
    FROM ventas v
    JOIN automoviles_stock a ON v.numero_serie = a.numero_serie
    JOIN concesionarios c ON a.id_concesionario = c.id_concesionario
    WHERE c.id_concesionario = 1
    LIMIT 5;
]])
system.print(q_limit)

print("\n--- B. VENTAS ORDENADAS MAYORES A $150,000.00 (ORDER BY) ---")
local q_order = db:query([[
    SELECT id_venta, numero_serie, precio_venta, modo_pago, fecha_entrega 
    FROM ventas 
    WHERE precio_venta > 150000.00 
    ORDER BY precio_venta DESC;
]])
system.print(q_order)

print("\n--- C. MODIFICAR NOMBRE DE CONCESIONARIOS CON REPLACE ---")
assert(db:query([[
    UPDATE concesionarios 
    SET nombre = CONCAT(nombre, '-nacional');
]]))
local q_replace = db:query("SELECT id_concesionario, nombre FROM concesionarios LIMIT 5;")
system.print(q_replace)
print("> Interpretación: Se actualizó el nombre de cada concesionaria concatenando el sufijo nacional.")

print("\n--- D. ELIMINAR MODELOS ANTERIORES AL AÑO 2020 (DELETE) ---")
assert(db:multi_query([[
    SET FOREIGN_KEY_CHECKS = 0;
    DELETE FROM ventas WHERE numero_serie IN (
        SELECT s.numero_serie FROM automoviles_stock s 
        JOIN modelos_autos m ON s.id_modelo = m.id_modelo 
        WHERE m.anho < 2020
    );
    DELETE FROM automoviles_stock WHERE id_modelo IN (
        SELECT id_modelo FROM modelos_autos WHERE anho < 2020
    );
    DELETE FROM modelos_autos WHERE anho < 2020;
    SET FOREIGN_KEY_CHECKS = 1;
]]))
assert(db:query("DELETE FROM modelos_autos WHERE anho < 2020;"))
local q_delete = db:query("SELECT id_modelo, marca, modelo, anho FROM modelos_autos LIMIT 5;")
system.print(q_delete)
print("> Interpretación: Se removieron del catálogo todos los vehículos con año de fabricación menor a 2020.")

db:close()
\end{lstlisting}

\newpage

\section{Explicación Detallada de los Pasos Realizados}

\begin{enumerate}
    \item \textbf{Diseño y Creación de Esquema DDL}: Se estructuró la base de datos \texttt{autos\_concesionario\_db} en MariaDB definiendo claves primarias, foráneas y restricciones \texttt{NOT NULL}.
    \item \textbf{Generación Dinámica e Inserción de Datos}: Mediante ciclos en LuaJIT se sembraron 20 registros por tabla principal y 25 simulaciones en \texttt{ventas\_acumuladas}.
    \item \textbf{Filtrado con LIMIT}: Se ejecutó una consulta \texttt{JOIN} filtrando ventas del concesionario 1 acotadas a 5 registros.
    \item \textbf{Ordenamiento con ORDER BY}: Se seleccionaron transacciones con precio mayor a \$150,000.00 ordenadas descendentemente.
    \item \textbf{Actualización Masiva (UPDATE/CONCAT)}: Se modificaron los nombres de las concesionarias añadiendo la extensión \texttt{-nacional}.
    \item \textbf{Eliminación Condicional (DELETE)}: Se depuró el catálogo removiendo vehículos y modelos con año de fabricación inferior a 2020.
\end{enumerate}

\end{document}